\section{Taxonomy và phân bố dữ liệu}
ROD Final giữ thứ tự 25 lớp từ ID 0 đến 24: Bike, Building, Car, Person, Stairs, Traffic sign, Electrical Pole, Road, Motorcycle, Dustbin, Dog, Manhole, Tree, Guard rail, Pedestrian crosswalk, Truck, Bus, Bench, Traffic Cone, Fire hydrant, Teraffic Barrel, Plant Pot, Electrical Box, Chair và Bicycle Rack. Chuỗi ``Teraffic Barrel'' được giữ đúng theo \texttt{data.yaml}. Nếu chỉ sửa tên ở một nơi, checkpoint và giao diện vẫn chạy nhưng nhãn hiển thị có thể lệch với evaluator; taxonomy vì vậy là một phần của hợp đồng dữ liệu.

Các lớp bao phủ ba nhóm chính. Nhóm tác nhân chuyển động gồm người, xe đạp, xe máy, ô tô, xe buýt, xe tải và chó. Nhóm cấu trúc đường đi gồm Road, Stairs, Manhole, Guard rail và Pedestrian crosswalk. Nhóm vật thể ven đường gồm biển báo, cột điện, thùng rác, ghế, chậu cây, hộp điện, cọc tiêu, thùng chắn và giá xe đạp. Phân nhóm chỉ hỗ trợ thảo luận; model vẫn học 25 lớp độc lập và không có nhãn cấp cha.

Phân bố lớp không cân bằng. mAP cho mỗi lớp trọng số bằng nhau, trong khi micro Precision/Recall bị chi phối nhiều hơn bởi lớp có nhiều object. Một mô hình có micro Recall cao vẫn có thể bỏ sót lớp hiếm; ngược lại mAP cao không cho biết số cảnh báo sai thực tế tại ngưỡng vận hành. Vì thế báo cáo đọc đồng thời chỉ số tổng, TP/FP/FN và AP theo lớp.

\begin{figure}[H]\centering
\includegraphics[width=0.88\textwidth]{fig_class_distribution.png}
\caption{Phân bố đối tượng theo lớp ở lần phân tích dữ liệu nguồn; số liệu Final lấy từ báo cáo đóng băng.}\label{fig:classdist}
\end{figure}

\section{Chuỗi làm sạch và đóng băng}
Bản nguồn có 24.326 ảnh và 40.195 object. V2 còn 24.187 ảnh và 40.041 object sau khi loại 139 cặp ảnh--nhãn cùng bốn annotation sai được xác nhận. V3 giữ nguyên quy mô và áp dụng bảy relabel Bike sang Motorcycle đã được duyệt; 152 đề xuất khác và 18 bbox chưa sửa. V4 cách ly 929 cặp có cạnh leakage trực tiếp, đưa dữ liệu về 23.258 ảnh và 38.797 object. Final cách ly thêm 178 cặp, gồm 146 TRAIN và 32 VALID, không thay đổi TEST.

\begin{table}[H]\centering
\caption{Thay đổi quy mô dữ liệu qua các mốc làm sạch}\label{tab:cleaning}
\begin{tabular}{lrrp{6.1cm}}\toprule
Mốc & Ảnh & Object & Thay đổi chính\\\midrule
Nguồn & 24.326 & 40.195 & Chuyển đổi nhãn và kiểm toán độc lập.\\
V2 & 24.187 & 40.041 & Loại 139 cặp và annotation sai đã xác nhận.\\
V3 & 24.187 & 40.041 & Relabel 7 trường hợp Bike sang Motorcycle.\\
V4 & 23.258 & 38.797 & Cách ly 929 cặp có leakage trực tiếp.\\
Final A4 & 23.080 & 38.467 & Cách ly thêm 178 cặp, giữ nguyên TEST.\\\bottomrule
\end{tabular}
\end{table}

Final còn sửa 178 orphan label trong staging sau khi đối chiếu checksum với manifest. Báo cáo xác nhận 0 ảnh thiếu label, 0 label thiếu ảnh, 0 lỗi class ID/cú pháp/hình học và 0 nhóm exact duplicate xuyên split. Tuy vậy còn 271 nhóm exact duplicate trong cùng split, 18 đề xuất bbox và các relabel Bike/Motorcycle chưa duyệt. Chúng là giới hạn dữ liệu, không phải lỗi đã xử lý.

\section{Nguyên tắc chia tập và chống rò rỉ}
TRAIN dùng cập nhật trọng số. VALID dùng theo dõi checkpoint, chọn ngưỡng và so sánh mô hình. TEST chỉ ước lượng cuối sau khi mọi quyết định đã khóa. Nếu ngưỡng được hạ vì một ảnh TEST khó, ảnh đó không còn độc lập và kết quả sẽ lạc quan. Bởi vậy script TEST kiểm hash dữ liệu, checkpoint, ngưỡng và protocol trước khi chạy.

Exact hash phát hiện bản sao byte giống hệt, còn manifest xử lý các quan hệ leakage đã xác định. Nó không chứng minh đã loại mọi near-duplicate do crop, resize hoặc nén lại. Kiểm toán tiếp theo có thể dùng embedding hay perceptual hash với duyệt thủ công; không nên tự động loại ảnh chỉ vì khoảng cách gần nếu chưa xác nhận cùng cảnh.
